\documentclass[11pt,a4paper]{article}


\usepackage[magyar]{babel}
\usepackage{amsmath,amssymb,amsthm}
\usepackage{hyperref}

\hypersetup{
  colorlinks=true,
  linkcolor=blue,
  urlcolor=blue,
  citecolor=blue
}

\newtheorem{tetel}{Tétel}[section]
\newtheorem{allitas}[tetel]{Állítás}
\newtheorem{lemma}[tetel]{Lemma}
\newtheorem{kovetkezmeny}[tetel]{Következmény}

\theoremstyle{definition}
\newtheorem{definicio}[tetel]{Definíció}
\newtheorem{pelda}[tetel]{Példa}
\newtheorem{megjegyzes}[tetel]{Megjegyzés}

\renewcommand{\proofname}{Bizonyítás}

\DeclareMathOperator{\Frac}{Frac}
\DeclareMathOperator{\Spec}{Spec}
\DeclareMathOperator{\ord}{ord}
\DeclareMathOperator{\length}{length}

\newcommand{\OO}{\mathcal O}
\newcommand{\mm}{\mathfrak m}
\newcommand{\pp}{\mathfrak p}
\newcommand{\qq}{\mathfrak q}
\newcommand{\ZZ}{\mathbb Z}
\newcommand{\NN}{\mathbb N}
\newcommand{\kk}{\Bbbk}

\title{Értékelések és értékelésgyűrűk\\
\large Különös tekintettel a diszkrét értékelésgyűrűkre}
\author{Előadásjegyzet az Algebrai Görbék c. tárgyhoz}
\date{}

\begin{document}

\maketitle
\tableofcontents

\section{Értékelések}

\subsection{Rendezett abel-csoportok}

\begin{definicio}
Egy $\Gamma$ abel-csoportot \emph{rendezett abel-csoportnak} nevezünk,
ha el van látva egy teljes rendezéssel, és a rendezés kompatibilis a
csoportművelettel: ha $\alpha\leq \beta$, akkor minden
$\gamma\in\Gamma$ esetén
\[
  \alpha+\gamma\leq \beta+\gamma.
\]
\end{definicio}





Az értékelések értékkészletét egy ilyen $\Gamma$ csoport
adja. A legfontosabb eset számunkra a $\Gamma=\ZZ$ ellátva a szokásos
rendezéssel. Ekkor beszélünk diszkrét értékelésről.

Formálisan hasznos a $\Gamma$ csoportot kiegészíteni egy $\infty$
elemmel, amelyre
\[
  \alpha < \infty,\qquad
  \alpha+\infty=\infty+\alpha=\infty
\]
minden $\alpha\in\Gamma$ esetén. Így kényelmesen definiálhatjuk
$
\nu(0)=\infty$ értéket.

\subsection{Értékelések definíciója}

\begin{definicio}
\label{def:ertekeles}
Legyen $K$ test, $\Gamma$ rendezett abel-csoport. Egy
\emph{értékelés} a $K$ testen egy olyan
\[
  \nu:K^\times\longrightarrow \Gamma
\]
leképezés, amelyre minden $x,y\in K^\times$ esetén
\[
  \nu(xy)=\nu(x)+\nu(y),
\]
továbbá ha $x+y\neq 0$, akkor
\[
  \nu(x+y)\geq \min\{\nu(x),\nu(y)\}.
\]
\begin{kovetkezmeny}
   Triviális értékelés: $\nu \equiv0 $
\end{kovetkezmeny}
\begin{allitas}

Legyen $\nu$ értékelés $K$-n. Ekkor:
\begin{enumerate}
  \item $\nu(1)=0$ és $\nu(-1)=0$;
  \item $\nu(x^{-1})=-\nu(x)$ minden $x\in K^\times$ esetén;
  \item ha $\nu(x)\neq \nu(y)$, akkor $\nu(x+y)=\min\{\nu(x),\nu(y)\}$.
\end{enumerate}
\end{allitas}

\begin{proof}
Az $1\cdot 1=1$ azonosságból
$
\nu(1)=
\nu(1)+
\nu(1)$, tehát $
\nu(1)=0$. Mivel $(-1)^2=1$,
$2v(-1)=0$, ezért $
\nu(-1)=0$ rendezett abel-csoportban, hiszen abban
nincs torzió. Az $xx^{-1}=1$ azonosságból adódik
$
\nu(x)+
\nu(x^{-1})=0$, tehát $
\nu(x^{-1})=-
\nu(x)$.

Tegyük fel, hogy $
\nu(x)<
\nu(y)$. A \ref{def:ertekeles}. definícióban
szereplő egyenlőtlenség szerint $
\nu(x+y)\geq 
\nu(x)$. Ha szigorú
egyenlőtlenség állna fenn, akkor
$
\nu(y)=
\nu((x+y)-x)\geq \min\{
\nu(x+y),
\nu(x)\}=
\nu(x),
$
ami ellentmond a kiinduló feltevésnek, következésképpen
$
\nu(x+y)=
\nu(x)$.
\end{proof}

\begin{definicio}
Az értékelés \emph{értékcsoportja} a $
\nu(K^\times)\subseteq\Gamma$
rendezett részcsoport. Az értékelés \emph{triviális}, ha
$
\nu(x)=0$ minden $x\in K^\times$ esetén.
\end{definicio}
Megállapodás szerint $
\nu(0)=\infty$.
\end{definicio}

Az első tulajdonság azt mondja, hogy $
\nu$ csoporthomomorfizmus
$K^\times$ multiplikatív csoportjából $\Gamma$ additív csoportjába. A
második tulajdonság az ultrametrikus háromszög-egyenlőtlenség
additív alakja.

\begin{allitas}
$\OO=\OO_\nu:={x\in F |\nu(x)\geq0}$ részgyűrű $F$-ben. $\OO_\nu^*=\{x\in F |\nu(x)=0\}$
\end{allitas}
\begin{proof}
$0\in \OO$, mert $\nu(0)= \infty>0$, $1\in \ker(\nu) \Rightarrow \nu(1)=0\geq0\Rightarrow 1\in \OO$
$x,y\in \OO $ esetén persze $x+y$, $xy\in\OO$

\end{proof}
\begin{definicio}
    $\OO_\nu$ a $\nu$ értékelés gyűrűje
\end{definicio}













\section{Értékelésgyűrűk}

\subsection{Az értékeléshez tartozó gyűrű}

\begin{definicio}
Legyen $
\nu$ értékelés a $K$ testen. Az
\[
  \OO_v=\{x\in K\mid 
\nu(x)\geq 0\}
\]
részgyűrűt a $
\nu$ értékelés \emph{értékelésgyűrűjének} nevezzük. Az
\[
  \mm_v=\{x\in K\mid 
\nu(x)>0\}
\]
ideált az \emph{értékeléshez tartozó maximális ideálnak} nevezzük.
\end{definicio}

\begin{allitas}
    $\OO_\nu^*= \{x\in \OO_\nu \}\bigcap\{x^{-1}\in \OO_\nu \}$
\end{allitas}


\begin{proof}
Mert  ekkor $\nu(x)\geq0$, és $\nu(x^{-1})\geq0 \Leftrightarrow -\nu(x)\geq0$, szóval $\nu(x) =0$
\end{proof}

\begin{megjegyzes}
    $\nu$ és $\OO_\nu$ megadása egyenértékű
\end{megjegyzes}

\begin{kovetkezmeny}
    Ha $\nu$ értékelés akkor $n\in \mathbb{N}$-re $n\cdot \nu$ is értékelés.
\end{kovetkezmeny}
\subsection{Ekvivalens értékelések}

\begin{definicio}
Legyen $\nu,\nu'$ két értékelés ugyanazon az $F$ testen.
Azt mondjuk, hogy $\nu$ és $\nu'$ \emph{ekvivalensek},
ha létezik olyan $n,m\in\mathbb Z_{>0}$, hogy
\[
  n\nu = m\nu'.
\]
\end{definicio}

\begin{megjegyzes}
Minden nem triviális értékelés ekvivalens egyetlen
szürjektív értékeléssel, ezt normalizált érékelésnek is hívjuk.
\end{megjegyzes}
\begin{definicio}
    $\pi \in F$ uniformizáló, ha $\nu(\pi)=1$
\end{definicio}
\begin{allitas}
    $\OO_\nu$-nek pontosan egy nem triviális maximális ideálja van, ez:
    $m:= \{x\in\OO_\nu|\nu(x)>0\}=(\pi)$.
\end{allitas}
$\forall I \lhd \OO_\nu$-re $\exists n: I=(\pi^n)$ $\Rightarrow$ $\OO_\nu$ főideálgyűrű.

\begin{megjegyzes}
    $\OO_\nu \leq F \Rightarrow$ integritási tartomány
\end{megjegyzes}
\begin{kovetkezmeny}
   $ \OO_\nu$ UFD, $\forall 0 \neq x \in \OO_\nu$ egyértelműen $x= u\cdot \pi^n$, $n=\nu(x)$
\end{kovetkezmeny}
\begin{definicio}
    Az ilyen integritási tarományok a diszkrét értékelés gyűrűk (DVR)
\end{definicio}

\begin{proof}
    Ha $x\in \OO_\nu$, akkor: \\
    $\nu(x)=0 \Rightarrow x\in \OO_\nu^*$ vagy \\
    $\nu(x)>0 \Rightarrow m\{x\in\OO_\nu|\nu(x)>0\}$ \\
    Ez ideál, mert: \\
    $\nu(0)=\infty \Rightarrow 0\in m$ \\
    $x\in m, y \in \OO \Rightarrow \nu (xy)= \nu(x+y)>0$ \\
    $x,y\in m \Rightarrow \nu(x+y)\geq \min(\nu(x),\nu(y))>0$ \\
    és persze egyetlen maximális, mert minden nem invertálhatót tartalmaz.
     
     
     Legyen $x\in \OO_\nu\backslash\{0\}$ \\
    $\rightarrow nu(x) = n\in \mathbb{Z}_{\geq0}$\\
    $\Rightarrow x = u\cdot \pi^n$, $u\in F$, hogy $\nu(u)=0$ $\Rightarrow u\in \OO_\nu^*$\\
    $I\rhd \OO_\nu \Rightarrow d = \min\{ \nu(x)|x\in I\}\in\mathbb{Z}_{\geq0}$ \\
    $\Rightarrow \exists x\in I \nu(x)=d\Rightarrow x=u\cdot\pi^d\Rightarrow \pi^d=u^{-1}x\in I\Rightarrow(\pi^d)\subseteq I$ \\ 
    $y\in I \rightarrow \nu(y) \geq d \Rightarrow y = v\cdot\pi^{\hat{d}}$ \\
    $\hat{d}\geq 0, v \in \OO_\nu^*$\\
    $\Rightarrow (\pi d)\supseteq I$
\end{proof}
\begin{kovetkezmeny}
    $x\in \OO_\nu$ pontosan akkor irred/prím, ha uniformáló
\end{kovetkezmeny}

\begin{tetel}
    Legyen $\OO$ egy DVR, $F$ az $\OO$ hányadosteste. Ekkor $\exists!\nu:F\rightarrow\mathbb{Z}\cup\{\infty\}$ értékelés úgy, hogy $\OO = \OO_\nu$
\end{tetel}
\begin{proof}
    $0 \neq \frac{u}{v}\in F,u,v\in\OO, v\neq0$ \\
    $\rightarrow v=v^*\pi^n$ $n\in\mathbb{N}$ \\
    $v^*\in\OO^*\Rightarrow\frac{u}{v}=\frac{u(v^*)^{-1}}{\pi^n}$\\
    $u(v^*)^{-1}=w\pi^m, m\in\mathbb{N}$, $ \nu(\frac{u}{v})=mn= \nu(u)-\nu(v)$
\end{proof}
\begin{allitas}
Az $\OO_v$ értékelésgyűrű lokális integritási tartomány, maximális ideálja
$\mm_v$, egységcsoportja pedig
\[
  \OO_v^\times=\{x\in K^\times\mid 
\nu(x)=0\}.
\]
Továbbá $\Frac(\OO_v)=K$.
\end{allitas}

\begin{proof}
Az értékelés tulajdonságaiból közvetlenül következik, hogy $\OO_v$
zárt az összeadásra és szorzásra, továbbá $1\in\OO_v$. Mivel
$\OO_v\subseteq K$, ezért $\OO_v$ is integritási tartomány.

Ha $x\in\OO_v$ és $
\nu(x)=0$, akkor $
\nu(x^{-1})=0$, tehát
$x^{-1}\in\OO_v$, vagyis $x$ egység. Ha $
\nu(x)>0$, akkor
$
\nu(x^{-1})<0$, tehát $x^{-1}\notin\OO_v$, azaz $x$ nem egység.
Így az egységek pontosan a
$0$ értékelésű elemek, a nemegységek pedig pontosan $\mm_v$ elemei. Egy
gyűrű pontosan akkor lokális, ha a nemegységek halmaza ideál; itt ez
az $\mm_v$ ideálra teljesül.

Végül ha $x\in K^\times$, akkor vagy $
\nu(x)\geq 0$, és ekkor
$x\in\OO_v$, vagy $
\nu(x)<0$, és ekkor $
\nu(x^{-1})>0$, tehát
$x^{-1}\in\OO_v$. Így minden $K$-beli elem két $\OO_v$-beli elem
hányadosaként írható, tehát $\Frac(\OO_v)=K$.
\end{proof}

\begin{tetel}[Értékelésgyűrűk intrinzikus jellemzése]
Legyen $K$ test, és legyen $R\subseteq K$ olyan részgyűrű, amelyre
$\Frac(R)=K$. Ekkor $R$ pontosan akkor valamely értékelés
értékelésgyűrűje, ha minden $x\in K^\times$ esetén
\[
  x\in R \quad\text{vagy}\quad x^{-1}\in R.
\]
\end{tetel}

\begin{proof}
Ha $R=\OO_v$, akkor az állítás az előző bizonyítás végén már
megjelent.

Fordítva tegyük fel, hogy $R$ rendelkezik a fenti tulajdonsággal.
Tekintsük a $K^\times/R^\times$ faktorcsoportot. Egyelőre
multiplikatív jelölést használunk, de rendezett abel-csoportként
később additív jelölésre térhetünk át. A rendezést definiáljuk úgy,
hogy
\[
  xR^\times\geq yR^\times
  \quad\Longleftrightarrow\quad
  xy^{-1}\in R.
\]
Ez jól definiált, mert egységgel való szorzás nem változtatja meg az
$R$-beli tagságot. A rendezés teljes, hiszen minden
$xy^{-1}\in K^\times$ elemre vagy $xy^{-1}\in R$, vagy
$(xy^{-1})^{-1}=yx^{-1}\in R$. Antiszimmetrikus is: ha
$xy^{-1}\in R$ és $yx^{-1}\in R$, akkor $xy^{-1}$ egység $R$-ben,
tehát $xR^\times=yR^\times$. A tranzitivitás és a csoportművelettel
való kompatibilitás közvetlenül következik abból, hogy $R$ gyűrű.

A

  $
\nu:K^\times\longrightarrow K^\times/R^\times,\qquad
  x\longmapsto xR^\times
$
leképezés csoporthomomorfizmus. Ha $x+y\neq 0$ és például
$
\nu(x)\geq 
\nu(y)$, akkor $xy^{-1}\in R$, ezért
\[
  (x+y)y^{-1}=xy^{-1}+1\in R.
\]
Ez azt jelenti, hogy $
\nu(x+y)\geq 
\nu(y)=\min\{
\nu(x),
\nu(y)\}$. A másik
eset ugyanilyen. Így $
\nu$ értékelés, és definíció szerint
\[
  \OO_v=\{x\in K\mid 
\nu(x)\geq 
\nu(1)\}=R.
\]
\end{proof}

\begin{megjegyzes}
A fenti tétel miatt az értékelésgyűrűt gyakran azzal a
tulajdonsággal is definiálják, hogy minden $x\in K^\times$ elemre
$x$ vagy $x^{-1}$ benne van a gyűrűben. Ez a megfogalmazás nagyon
hasznos, amikor nem egy konkrét képlettel adott értékelésből indulunk
ki, hanem egy lokális gyűrű tulajdonságait vizsgáljuk.
\end{megjegyzes}

\subsection{Alapvető tulajdonságok}

\begin{allitas}
Egy értékelésgyűrű ideáljai lineárisan rendezettek a tartalmazásra
nézve: ha $I,J\subseteq\OO_v$ ideálok, akkor $I\subseteq J$ vagy
$J\subseteq I$.
\end{allitas}

\begin{proof}
Tegyük fel, hogy $I\not\subseteq J$. Válasszunk $a\in I\setminus J$
elemet. Megmutatjuk, hogy $J\subseteq I$. Legyen $b\in J$. Ha
$b=0$, nincs mit bizonyítani. Ha $b\neq 0$, akkor $a/b\notin\OO_v$,
mert ellenkező esetben $a=(a/b)b\in J$ lenne. Értékelésgyűrűről
lévén szó, ebből $b/a\in\OO_v$ következik. Így
$b=(b/a)a\in I$, tehát $J\subseteq I$.
\end{proof}

\begin{kovetkezmeny}
Egy értékelésgyűrű prímideáljai lineárisan rendezettek. Konkrétan
egy értékelésgyűrű spektruma lánc.
\end{kovetkezmeny}

\begin{allitas}
Minden értékelésgyűrű egészzárt a hányadostestében.
\end{allitas}

\begin{proof}
Legyen $R=\OO_v$, $K=\Frac(R)$, és tegyük fel, hogy
$x\in K$ egész $R$ fölött. Ha $x\notin R$, akkor
$x^{-1}\in R$ és $
\nu(x)<0$. Mivel $x$ egész $R$ fölött, teljesít egy
olyan egyenletet, hogy
\[
  x^n+a_{n-1}x^{n-1}+\cdots+a_1x+a_0=0,
  \qquad a_i\in R.
\]
Osszunk $x^{n-1}$-gyel:
\[
  x=-(a_{n-1}+a_{n-2}x^{-1}+\cdots+a_0x^{-(n-1)}).
\]
A jobb oldal $R$-ben van, hiszen $a_i\in R$ és $x^{-1}\in R$. Ez
ellentmond $x\notin R$-nek. Tehát $x\in R$, vagyis $R$ egészzárt.
\end{proof}

\section{Diszkrét értékelések}

\subsection{Definíció és normalizálás}

\begin{definicio}
Egy $
\nu:K^\times\to\Gamma$ értékelés \emph{diszkrét értékelés}, ha
nem triviális, és értékcsoportja rendezett csoportként izomorf
$\ZZ$-vel. Ilyenkor a választott izomorfizmus után úgy tekintjük, hogy
\[
  
\nu:K^\times\longrightarrow \ZZ.
\]
Ha $
\nu(K^\times)=\ZZ$, akkor $
\nu$-t \emph{normalizált} diszkrét
értékelésnek nevezzük.
\end{definicio}

Ha az értékcsoport $\ZZ$ egy $e\ZZ$ részcsoportja, akkor az értékelést
$e$-vel osztva normalizálhatjuk. A gyakorlatban diszkrét értékelés
alatt legtöbbször normalizált diszkrét értékelést értünk.

\begin{definicio}
Legyen $
\nu$ normalizált diszkrét értékelés $K$-n, értékelésgyűrűje
$R=\OO_v$, maximális ideálja $\mm=\mm_v$. Egy $\pi\in R$ elemet
\emph{uniformizálónak} vagy \emph{lokális paraméternek} nevezünk, ha
\[
  
\nu(\pi)=1.
\]
\end{definicio}

\begin{allitas}
Legyen $R=\OO_v$ a $
\nu$ normalizált diszkrét értékeléshez tartozó
értékelésgyűrű, és legyen
$\pi$ uniformizáló. Ekkor minden $x\in K^\times$ egyértelműen írható
\[
  x=\nu\pi^n
\]
alakban, ahol $\nu\in R^\times$ és $n\in\ZZ$. Továbbá $x\in R$ pontosan
akkor, ha $n\geq 0$, és $x\in\mm$ pontosan akkor, ha $n>0$.
\end{allitas}

\begin{proof}
Legyen $n=
\nu(x)$. Ekkor
\[
  
\nu(x\pi^{-n})=
\nu(x)-n 
\nu(\pi)=0,
\]
tehát $\nu=x\pi^{-n}\in R^\times$, és $x=\nu\pi^n$. Az egyértelműség
abból következik, hogy ha $\nu\pi^n=\nu'\pi^{n'}$, akkor az értékelést
véve $n=n'$, majd $\nu=x\pi^{-n}=\nu'$. A tagsági állítások közvetlenül következnek
az értékelésgyűrű és a maximális ideál definíciójából.
\end{proof}

\begin{kovetkezmeny}
Ha $R$ diszkrét értékelésgyűrű és $\pi$ uniformizáló, akkor
\[
  \mm=(\pi),
  \qquad
  \mm^n=(\pi^n)\quad (n\geq 1).
\]
\end{kovetkezmeny}

\begin{proof}
Az $\mm$ elemei pontosan azok az $\nu\pi^n$ alakú elemek, ahol
$n\geq 1$, tehát mind oszthatók $\pi$-vel. Mivel $\pi\in\mm$,
$\mm=(\pi)$. A hatványokra vonatkozó állítás ebből következik.
\end{proof}

\subsection{A DVR mint lokális főideálgyűrű}

\begin{definicio}
Egy $R$ integritási tartományt \emph{diszkrét értékelésgyűrűnek},
röviden \emph{DVR-nek}, nevezünk, ha $R$ egy normalizált diszkrét
értékeléshez tartozó értékelésgyűrű.
\end{definicio}

\begin{tetel}[A DVR-ek ideáljai]
Legyen $R$ DVR, $\pi$ uniformizáló, $\mm=(\pi)$ maximális ideál.
Ekkor $R$ minden nem nulla ideálja $\mm^n=(\pi^n)$ alakú egy
egyértelmű $n\geq 0$ egészre. Következésképpen $R$ noether lokális
főideálgyűrű.
\end{tetel}

\begin{proof}
Legyen $I\neq 0$ ideál. Tekintsük az
\[
  S=\{
\nu(x)\mid 0\neq x\in I\}\subseteq \NN
\]
halmazt. Ez nem üres, ezért van legkisebb eleme, mondjuk $n$. Válasszunk
$a\in I$ elemet $
\nu(a)=n$ értékkel. Ekkor $a=\nu\pi^n$ valamely
$\nu\in R^\times$ egységgel, tehát $(a)=(\pi^n)$.

Mivel $a\in I$, kapjuk, hogy $(\pi^n)\subseteq I$. Másrészt minden
$x\in I$ esetén $
\nu(x)\geq n$, tehát $x$ osztható $\pi^n$-nel, vagyis
$x\in(\pi^n)$. Így $I=(\pi^n)$. Az $n$ egyértelmű, mert
$(\pi^n)=(\pi^m)$ esetén az ideálok nem nulla elemeinek minimális
értéke egyszerre $n$ és $m$.
\end{proof}

\begin{kovetkezmeny}
Egy DVR prímideáljai pontosan $(0)$ és $\mm$. Következésképpen $\dim R=1$.
\end{kovetkezmeny}

\begin{proof}
Legyen $\pp\neq (0)$ prímideál. Az előző tétel szerint
$\pp=(\pi^n)$ valamely $n\geq 1$ esetén. Mivel $\pi^n\in\pp$ és
$\pp$ prímideál, $\pi\in\pp$, tehát $\mm=(\pi)\subseteq\pp$. Mivel
$\mm$ maximális, $\pp=\mm$.
\end{proof}

\begin{kovetkezmeny}
Egy DVR egészzárt, noether, egydimenziós lokális integritási
tartomány.
\end{kovetkezmeny}

\begin{proof}
Az egészzártság minden értékelésgyűrűre igaz. A noetherség és
az egydimenziósság az előző tételekből következik.
\end{proof}

\subsection{Ekvivalens jellemzések}

\begin{tetel}[DVR-ek karakterizációja]
Legyen $R$ noether lokális integritási tartomány, amely nem test, és
legyen $\mm$ a maximális ideálja. Ekkor az alábbiak ekvivalensek:
\begin{enumerate}
  \item $R$ diszkrét értékelésgyűrű;
  \item $R$ lokális főideálgyűrű;
  \item $\mm$ főideál, és minden $0\neq x\in R$ egyértelműen
        $x=\nu\pi^n$ alakú, ahol $\nu\in R^\times$, $n\geq 0$;
  \item $R$ egydimenziós és egészzárt.
\end{enumerate}
\end{tetel}

\begin{proof}
Az $(1)\Rightarrow(2)$ implikációt az ideálokról szóló tétel adja.
Az $(2)\Rightarrow(3)$ világos: ha $R$ lokális főideálgyűrű és nem
test, akkor a maximális ideál $\mm=(\pi)$ valamely nem nulla
nemegységre. Ha $0\neq x\in R$, akkor $x$ valamely legnagyobb
$n\geq 0$ kitevőre $\mm^n$-ben van. Ilyen $n$ létezik, mert
$\bigcap_{n\geq 0}\mm^n=0$ a Krull-féle metszettétel miatt. Ekkor
$x=\pi^n \nu$ és $\nu\notin\mm$, tehát $\nu$ egység.

A $(3)\Rightarrow(1)$ implikációhoz legyen $K=\Frac(R)$. Ha
$0\neq z\in K$, írjuk $z=x/y$ alakban $x,y\in R\setminus\{0\}$. A
feltétel miatt
\[
  x=\nu\pi^m,\qquad y=w\pi^n
\]
egységekkel, ezért
\[
  z=(uw^{-1})\pi^{m-n}.
\]
Definiáljuk $
\nu(z)=m-n$. Az egyértelmű felbontás miatt ez jól
definiált, és azonnal teljesíti
$
\nu(zz')=
\nu(z)+
\nu(z')$, valamint
$
\nu(z+z')\geq\min\{
\nu(z),
\nu(z')\}$. Ekkor
\[
  R=\{z\in K\mid 
\nu(z)\geq 0\},
\]
tehát $R$ DVR.

Az $(1)\Rightarrow(4)$ már következményként szerepelt.

Végül bizonyítjuk a $(4)\Rightarrow(3)$ irányt. Ez a karakterizáció
legfontosabb része. Válasszunk $0\neq a\in\mm$ elemet. Mivel
$R$ egydimenziós noether lokális integritási tartomány, az $R/(a)$ gyűrű
lokális artin, ezért létezik $n\geq 1$ úgy, hogy
\[
  \mm^n\subseteq (a).
\]
Válasszuk $n$-t minimálisnak, és válasszunk
$b\in\mm^{n-1}$ elemet úgy, hogy $b\notin(a)$. Ekkor
$\mm b\subseteq(a)$, tehát ha $x=b/a\in K$, akkor $x\notin R$, de
$x\mm\subseteq R$.

Megmutatjuk, hogy $x\mm=R$. Ha $x\mm\subseteq\mm$ lenne, akkor a
végesen generált $\mm$ ideál stabil lenne az $x$-szel való szorzásra.
A determinánstrükk szerint ekkor $x$ egész lenne $R$ fölött.
Mivel $R$ egészzárt, ebből $x\in R$ következne, ellentétben
$b\notin(a)$ választásával. Tehát $x\mm\nsubseteq\mm$, és mivel $R$
lokális, $x\mm$ tartalmaz egységet. Így $x\mm=R$.

Ekkor létezik $\pi\in\mm$, amelyre $x\pi=1$. Ha $y\in\mm$, akkor
$xy\in R$, és
\[
  y=(xy)\pi\in(\pi).
\]
Tehát $\mm=(\pi)$ főideál. A felbontás
$0\neq r=\nu\pi^m$ alakban ismét a Krull-féle metszettételből adódik:
van legnagyobb $m$ úgy, hogy $r\in\mm^m$, és ekkor
$r\in(\pi^m)$, de $r\notin(\pi^{m+1})$, tehát $r=\nu\pi^m$ egy
$\nu\notin\mm$ egységgel.
\end{proof}

\begin{megjegyzes}
A tétel görbékhez leginkább használt alakja a következő: egy
egydimenziós noether lokális integritási tartomány pontosan akkor DVR, ha
egészzárt. Másképpen: normális görbén a zárt pontok lokális
gyűrűi DVR-ek.
\end{megjegyzes}

\section{Számolás DVR-ben}

\subsection{Rend, zérus és pólus}

Legyen $R$ DVR hányadosteste $K$, uniformizálója $\pi$, normalizált
értékelése $
\nu$. Egy $f\in K^\times$ elemre
\[
  f=\nu\pi^n,\qquad \nu\in R^\times,\quad n\in\ZZ.
\]
Az $n=
\nu(f)$ egész számot az $f$ \emph{rendjének} nevezzük.

\begin{itemize}
  \item Ha $
\nu(f)>0$, akkor $f$ zérusa van az adott helyen, rendje
        $
\nu(f)$.
  \item Ha $
\nu(f)=0$, akkor $f$ reguláris és nem tűnik el az adott
        helyen.
  \item Ha $
\nu(f)<0$, akkor $f$ pólusa van az adott helyen, pólusrendje
        $-
\nu(f)$.
\end{itemize}

\begin{allitas}
Legyen $R$ DVR, $\pi$ uniformizáló, és $0\neq f\in R$. Ekkor
\[
  \length_R(R/(f))=
\nu(f).
\]
\end{allitas}

\begin{proof}
Írjuk $f=\nu\pi^n$ alakban, ahol $\nu$ egység. Ekkor
$(f)=(\pi^n)$, ezért elég $R/(\pi^n)$ hosszát kiszámítani. A lánc
\[
  R\supset(\pi)\supset(\pi^2)\supset\cdots\supset(\pi^n)\supset 0
\]
$R/(\pi^n)$ kompozíciósorát adja, és az egymást követő faktorok
mind izomorfak $R/(\pi)$-vel. Így a hossz $n=
\nu(f)$.
\end{proof}

\subsection{Példák}

\begin{pelda}[$p$-adikus rend a racionális számokon]
Legyen $p$ prímszám. Minden $0\neq x\in\mathbb Q$ egyértelműen írható
\[
  x=p^n\frac{a}{b},
\]
ahol $n\in\ZZ$, $a,b\in\ZZ$, és $p\nmid a,b$. Ekkor
$v_p(x)=n$ diszkrét értékelés. Értékelésgyűrűje
\[
  \OO_{v_p}=\left\{\frac{a}{b}\in\mathbb Q\mid p\nmid b\right\}
  =\ZZ_{(p)},
\]
maximális ideálja pedig $p\ZZ_{(p)}$.
\end{pelda}

\begin{pelda}[$t$-adikus rend]
Legyen $k$ test. A $k(t)$ racionális függvénytestben a $t=0$ pontbeli
rend
\[
  v_t\left(t^n\frac{f(t)}{g(t)}\right)=n,
  \qquad f(0)g(0)\neq 0
\]
diszkrét értékelés. Értékelésgyűrűje $k[t]_{(t)}$, maximális ideálja
$(t)$.
\end{pelda}

\begin{pelda}[Formális hatványsorok]\leavevmode\par
A $k[[t]]$ formális hatványsorgyűrű DVR. Uniformizálója $t$, és
minden nem nulla hatványsor egyértelműen
\[
  a_n t^n+a_{n+1}t^{n+1}+\cdots
  =t^n(a_n+a_{n+1}t+\cdots)
\]
alakú, ahol $a_n\neq 0$, a zárójelben álló hatványsor pedig egység.
A hányadostest $k((t))$, a Laurent-sorok teste.
\end{pelda}

\end{document}
